\documentclass[12pt,a4paper]{article}
\usepackage{fontspec}
% Prefer installed fonts; bundled unmodified fonts keep CI and local builds portable.
\IfFontExistsTF{Liberation Sans}{\setmainfont{Liberation Sans}}{%
  \setmainfont{LiberationSans}[Path=fonts/,Extension=.ttf,UprightFont=*-Regular,BoldFont=*-Bold,ItalicFont=*-Italic,BoldItalicFont=*-BoldItalic]}

\usepackage{babel}
\babelprovide[main,import]{portuguese}
\usepackage{hyperref,booktabs,array,xcolor}
\usepackage[authoryear,round]{natbib}
\usepackage{xurl}
\hypersetup{citecolor=black}
\hypersetup{colorlinks=true,urlcolor=blue,linkcolor=black}
\usepackage[left=3cm,top=3cm,right=2cm,bottom=2cm]{geometry}
\usepackage{setspace}\onehalfspacing
\setcounter{tocdepth}{3}\setcounter{secnumdepth}{3}
\begin{document}\sloppy
\begin{titlepage}\centering CETI LUCAS MEIRELES ALVES\\CURSO TÉCNICO EM ANÁLISE E DESENVOLVIMENTO DE SISTEMAS\\PROJETO INTEGRADOR\vfill
{\Large\bfseries ESTUDAOFFLINE: PLANEJAMENTO DE ESTUDOS COM DADOS LOCAIS}\vfill
Teresina — Chapadinha Sul\\Outubro de 2026\end{titlepage}
\begin{titlepage}\centering Estudantes — identificação omitida nesta demonstração\\Professor: Aloisio\vfill
{\Large\bfseries ESTUDAOFFLINE: PLANEJAMENTO DE ESTUDOS COM DADOS LOCAIS}\vspace{2cm}
\begin{flushright}\begin{minipage}{.6\textwidth}Demonstração técnica de Projeto Integrador do curso técnico em Análise e Desenvolvimento de Sistemas do CETI Lucas Meireles Alves. Revisão e avaliação pelo professor pendentes.\end{minipage}\end{flushright}\vfill Teresina — Chapadinha Sul\\Outubro de 2026\end{titlepage}
\renewcommand{\contentsname}{Sumário}\tableofcontents\newpage
\section{Introdução}O EstudaOffline é um produto mínimo viável (MVP) de planejamento de estudos em Flutter para a web, com armazenamento local e proposta de funcionamento como aplicação web progressiva (PWA). O problema técnico é manter tarefas acessíveis no navegador quando a conexão está indisponível, reconhecendo limites de cache, armazenamento e atualização. O contexto escolar é o CETI Lucas Meireles Alves, curso técnico em Análise e Desenvolvimento de Sistemas, Teresina, bairro Chapadinha Sul, conforme o contexto registrado no protocolo \citep{projeto}.
\subsection{Objetivos}\subsubsection{Objetivo geral}Desenvolver e verificar tecnicamente um MVP de planejamento de estudos com persistência local em Flutter web.
\subsubsection{Objetivos específicos}\begin{enumerate}\item Implementar o cadastro e o acompanhamento de tarefas com exportação e importação de backup.\item Verificar persistência, validação de dados e carregamento offline em cenários técnicos documentados.\end{enumerate}
\subsection{Metodologia}O planejamento de avaliação considera a distinção entre episódios formativos e somativos e entre condições artificiais e naturalistas apresentada no FEDS. Este estudo delimita uma avaliação técnica artificial; isso não substitui validação com estudantes nem certifica o artefato \citep{eng17}.

O estudo é aplicado e técnico, baseado em implementação e cenários reproduzíveis. Utilizam-se exclusivamente tarefas sintéticas, sem participantes humanos, contas ou coleta de dados de usuários reais. Cada cenário deve registrar ambiente, versão, comando ou passos, resultado esperado, observado e artefato de evidência. A revisão estática e os testes automatizados examinam regras de tarefa, persistência e backup. A verificação no navegador deve distinguir persistência de dados e carregamento da aplicação offline: salvar tarefas não demonstra que todos os recursos da interface estejam em cache, distinção estabelecida no protocolo técnico \citep{projeto,sw}.

Os cenários contemplam criação e conclusão de tarefas, recarga, exportação e importação de backup, rejeição de dados inválidos e reabertura offline após instalação do cache. Resultados são restritos ao ambiente executado. Uma falha deve permanecer registrada; um cenário sem execução não pode ser contado como aprovado, conforme os critérios do protocolo \citep{projeto}.
\clearpage
\section{Referencial teórico}\label{sec:referencial-inicio}
\subsection{Organização conceitual e alcance da fundamentação}
O referencial do EstudaOffline articula seis eixos: organização do estudo, mediação docente, acesso, interação, proteção dos registros e engenharia verificável; essa organização é uma síntese proposta para o projeto, apoiada em revisões educacionais e de engenharia, e não uma teoria nova já validada \citep{education_palalas2020,eng04}.
A pergunta orientadora permanece técnica: em condições documentadas, o aplicativo permite organizar tarefas e preservar registros sintéticos quando a conexão é interrompida; a literatura oferece justificativas e limites para essa pergunta, mas não substitui a execução dos cenários previstos \citep{eng06}.

Revisões sistemáticas, mapeamentos e revisões guarda-chuva constituem fontes prioritárias porque permitem examinar conjuntos de pesquisas, em vez de tratar um resultado isolado como universal; entretanto, seus desenhos, populações e formas de síntese precisam ser considerados antes da transferência de conclusões \citep{education_prasse2024,eng04}.
Neste texto, convergência entre publicações significa coerência conceitual para uma decisão de projeto, e não uma nova estimativa estatística de efeito, pois as revisões podem compartilhar estudos primários e investigar intervenções diferentes \citep{education_prasse2024,education_palalas2020}.


A seleção desta fundamentação é narrativa e intencional: procura conceitos adequados ao problema técnico, sem estimar efeito educacional nem demonstrar cobertura exaustiva do campo \citep{SU11}.
A conferência dirigida passou a registrar, para os enunciados centrais, a localização consultada, o desenho da fonte, a sustentação disponível e os limites de transferência, mantendo distintos texto consultado, interpretação e proposta do projeto \citep{SU11}.

\subsection{Aprendizagem autorregulada e planejamento de tarefas}
Adota-se como referência organizadora a descrição do modelo cíclico de Zimmerman apresentada na revisão de Panadero, que distingue antecipação, execução e autorreflexão e inclui processos cognitivos, motivacionais e afetivos \citep{SU10}.
Essa escolha explicita a lente teórica, sem pressupor implementação de todos os processos: os registros atuais de tarefa, prazo e duração planejada são recursos organizacionais, enquanto a marcação de conclusão é apenas um estado declarado \citep{SU10}.
A revisão de Wolters e Brady situa gestão do tempo no quadro autorregulatório, mas não transforma esses campos em medidas validadas de motivação ou aprendizagem \citep{education_wolters2021}.

Palalas e Wark descrevem uma relação dinâmica entre aprendizagem móvel e autorregulação e recomendam integração curricular intencional, mas também relatam resultados neutros ou desfavoráveis quando se examina o uso de dispositivos \citep{education_palalas2020}.
As limitações da revisão incluem seleção em inglês, contextos formais e lacunas de relato em estudos incluídos; esses pontos impedem converter a recomendação geral em efeito esperado para este MVP \citep{education_palalas2020}.
A integração das tarefas à orientação do professor é uma proposta de uso a avaliar, e não um resultado produzido pelos cenários sintéticos atuais .

A revisão guarda-chuva de Prasse e colaboradores analisou 27 revisões sistemáticas e quatro meta-análises, identificando a importância de apoiar o ciclo autorregulatório de maneira integrada, com atenção às características dos suportes oferecidos \citep{education_prasse2024}.
No projeto, essa conclusão motiva examinar recursos organizacionais que poderiam apoiar parte do ciclo, sem afirmar que o usuário efetivamente planejou, monitorou ou refletiu ao utilizar a interface \citep{education_prasse2024}.
A revisão guarda-chuva inclui a revisão de Araka em seu conjunto analisado; por isso, citá-las juntas não representa duas bases empíricas independentes nem permite somar suas amostras \citep{education_prasse2024}.

Como operacionalização proposta, o cadastro de uma tarefa representa um registro de intenção, a mudança de estado representa uma declaração de andamento e a consulta posterior permite recuperar esse registro; esses eventos são observáveis no software, mas não medem diretamente compreensão, esforço cognitivo ou domínio de conteúdo \citep{education_araka2020}.
O critério verificável desta etapa é a correspondência entre os valores informados e recuperados, com identificação da versão e do cenário, enquanto eventuais indicadores educacionais dependerão de instrumentos e dados adicionais \citep{education_araka2020,eng11}.

Araka e colaboradores focalizam a medição e promoção da autorregulação em ambientes digitais de ensino superior; a transferência de instrumentos para estudantes do curso técnico não está demonstrada pelo simples uso da revisão \citep{education_araka2020}.
Nessa literatura, medidas ``offline'' são obtidas antes ou depois de um episódio de aprendizagem, ao contrário das medidas tomadas durante o processo; esse uso temporal do termo não demonstra funcionamento de software sem internet \citep{education_araka2020}.
No EstudaOffline, disponibilidade sem rede e preservação de registros são requisitos técnicos separados de qualquer inferência sobre autorregulação \citep{sw}.



\subsection{Desigualdade digital e continuidade em baixa conectividade}
A revisão de Martin, Ceviker e Gezer examina 77 estudos sobre desigualdade digital educacional e adota uma perspectiva multidimensional, envolvendo fatores, dimensões e intervenções, em vez de limitar o problema à presença de conexão \citep{martin2024}.
Para o EstudaOffline, essa perspectiva sustenta tratar continuidade sem rede como uma redução potencial de uma barreira técnica específica, reconhecendo que disponibilidade de dispositivos, competências e condições de uso continuam relevantes \citep{martin2024}.

A opção por manter registros no dispositivo é, portanto, uma decisão de escopo para experimentar continuidade de tarefas sob interrupção da rede; não constitui diagnóstico empírico da conectividade dos estudantes de Chapadinha Sul nem comprova inclusão digital na comunidade \citep{martin2024}.
O contexto escolar delimita o destino pretendido, mas a adequação local deverá ser investigada posteriormente, sem transferir automaticamente resultados de populações ou instituições analisadas pelas revisões \citep{martin2024,education_palalas2020}.

A continuidade operacional precisa distinguir disponibilidade da interface de disponibilidade dos registros, porque o armazenamento das tarefas e o cache dos recursos executáveis são mecanismos diferentes nas aplicações web \citep{local,sw}.
Por isso, o teste proposto combina preparação online, confirmação do controle pelo service worker, interrupção da rede e recarga; uma primeira visita sem recursos previamente armazenados constitui cenário distinto e deve ter seu limite explicitado \citep{sw}.

A escolha de Flutter corresponde à implementação do artefato, enquanto a justificativa científica se refere às funções e condições que serão examinadas; a documentação do framework não demonstra superioridade pedagógica nem garante, sozinha, operação offline \citep{flutter,education_palalas2020}.
O alvo inicialmente verificado é Flutter web, e extensões para Android ou compilação WebAssembly deverão registrar configuração, compatibilidade e resultados próprios, sem extrapolar os resultados de uma plataforma para outra \citep{eng04}.

\subsection{Acessibilidade, usabilidade e limites da avaliação técnica}
A revisão de Bong e Chen investiga a competência docente para produzir materiais e ambientes digitais acessíveis e inclusivos, trazendo a formação dos responsáveis pelo ensino para a discussão de acessibilidade \citep{bong2021}.
Essa perspectiva justifica incluir orientação de uso e clareza dos procedimentos no planejamento do trabalho, sem atribuir exclusivamente à interface a responsabilidade por remover todas as barreiras à participação \citep{bong2021}.

A revisão de Kumar e Mohite trata especificamente da usabilidade de aplicações de aprendizagem móvel, sustentando a necessidade de avaliação explícita das condições de interação \citep{kumar2017}.
No EstudaOffline, a proposta é examinar separadamente a execução correta das operações e a possibilidade de encontrar e compreender seus controles, pois uma função correta pode continuar difícil de utilizar \citep{kumar2017}.

O mapeamento de Kumar, Chand e Goundar reúne trabalhos sobre testes de usabilidade de aprendizagem móvel e identifica práticas como observação, questionários e avaliação com usuários, cujo emprego depende do contexto e do desenho da investigação \citep{kumarmapping2024}.
Como o estudo atual não possui participantes, não poderá oferecer estimativas de satisfação, taxa de sucesso dos estudantes ou representatividade de necessidades; inspeções técnicas preliminares serão relatadas como tais \citep{kumarmapping2024}.

A revisão de Gobbo e Bezerra discute acessibilidade digital na educação e amplia a atenção para barreiras que não se esgotam em aspectos técnicos; assim, corrigir um componente não autoriza concluir que a experiência educacional se tornou inclusiva \citep{gobbo2023}.
Os procedimentos de inspeção propostos passam a remeter a critérios identificáveis da WCAG 2.2: teclado (2.1.1), foco visível (2.4.7), identificação de erros (3.3.1), rótulos ou instruções (3.3.2), contraste (1.4.3) e nome, função e valor (4.1.2) \citep{wcag}.
Essa seleção parcial organiza verificações futuras e não comprova conformidade integral com a norma nem atendimento a todas as necessidades dos usuários \citep{wcag}.

Esses critérios organizam uma etapa de avaliação e não representam certificação obtida: a evidência deverá registrar controle examinado, procedimento, condição e resultado, inclusive falhas ou verificações ainda não realizadas \citep{eng06}.
A relação entre disponibilidade offline e acessibilidade é complementar, pois conservar a aplicação disponível não assegura que uma pessoa consiga operá-la; por isso, ambas as dimensões exigem evidências próprias \citep{martin2024,kumar2017}.

\subsection{Segurança, privacidade e preservação dos registros}
A revisão de Mkpojiogu, Hussain e Agbudu selecionou 21 estudos sobre segurança em aplicativos educacionais móveis e destacou limitações da literatura e dos enquadramentos utilizados, justificando examinar riscos sem pressupor que o contexto educacional torne uma aplicação segura \citep{OS18}.
A implementação atual não oferece contas nem telemetria própria, e os cenários de pesquisa usam dados sintéticos; porém, campos de texto livre podem receber informações pessoais se uma pessoa as inserir, o que impede prometer ausência universal desse conteúdo .
Registros locais na mesma origem e perfil do navegador não são isolados por identidade, e o uso de dispositivo compartilhado permanece uma ameaça específica a examinar \citep{local}.

A retenção no navegador depende das políticas de armazenamento e pode ser afetada por limpeza, cotas e condições de execução; portanto, persistência local deve ser entendida como comportamento condicionado e não como garantia de conservação permanente \citep{local,quota}.
O backup é proposto como caminho de recuperação, exigindo verificação de formato, versão, identificadores e conteúdo antes da substituição dos registros existentes, conforme os critérios definidos no protocolo \citep{eng06}.

Estudos sobre ataques com service workers e caches demonstram que mecanismos utilizados para disponibilidade também podem participar de ameaças de segurança, o que impede equiparar funcionamento offline a proteção dos dados \citep{OS01,OS02}.
As implicações propostas são delimitar recursos e escopo do cache, registrar atualização de versão e testar cenários de falha, sem afirmar vulnerabilidade encontrada nem proteção comprovada no EstudaOffline apenas pela leitura desses estudos \citep{OS01,OS02}.

O critério operacional de integridade consiste em recuperar os mesmos registros válidos e preservar o estado anterior quando uma importação inválida for rejeitada; trata-se de um requisito local verificável, e não de uma definição completa de segurança \citep{eng06}.
Essa delimitação evita interpretar aprovação em um teste de backup como evidência de resistência a todos os ataques, mantendo futuras análises de ameaças e testes de segurança como tarefas específicas \citep{OS18,eng06}.

\subsection{Qualidade do software, rastreabilidade e reprodutibilidade}
O estudo terciário de Zein, Salleh e Grundy organiza 24 estudos secundários de engenharia de aplicações móveis e evidencia a presença de diferentes fases e problemas de desenvolvimento, sustentando uma visão de qualidade que ultrapassa a construção da interface \citep{eng04}.
No projeto, isso justifica relacionar requisito, decisão de implementação, cenário e evidência, em vez de apresentar o aplicativo executável como conclusão suficiente da investigação \citep{eng04}.

O mapeamento de Júnior e colaboradores caracteriza técnicas de teste de requisitos não funcionais em 56 estudos e distingue atributos como confiabilidade, segurança e usabilidade, cujas avaliações demandam estratégias específicas \citep{eng06}.
Consequentemente, os testes de criação e conclusão de tarefas examinam operações funcionais, enquanto recarga sem rede, tratamento de falhas e interação assistiva exigem cenários próprios; uma dimensão aprovada não aprova automaticamente as demais \citep{eng06}.

A revisão de Soares e colaboradores sintetiza 101 estudos empíricos sobre integração contínua e identifica tanto efeitos positivos quanto desafios técnicos e de processo, impedindo tratar automação como ausência de defeitos \citep{eng15}.
No EstudaOffline, a integração contínua tem a função proposta de executar verificações, compilar artefatos e preservar registros associados à versão, enquanto a suficiência dos testes depende dos requisitos efetivamente examinados \citep{eng15}.

O estudo secundário de Anchundia e Fonseca identifica recursos de reprodução experimental cuja adequação depende do domínio; disponibilizar código é relevante, mas não substitui condições de execução e registros necessários à conferência \citep{eng11}.
Assim, cada cenário deve registrar versão, ambiente, dados, procedimento, resultado esperado e observado, permitindo conferir a observação sem presumir repetição idêntica em todos os navegadores \citep{eng11}.
O teste de interface atual utiliza armazenamento simulado em memória; sua reconstrução de widget não é uma recarga real do navegador, e os minutos registrados são duração planejada, não tempo efetivamente estudado .
Testes de conversão de backup sustentam regras do modelo, mas não demonstram sozinhos download real, atomicidade da restauração pela interface ou funcionamento offline \citep{eng06}.

\subsection{Modelo orientador e critérios de interpretação}
O encadeamento conceitual proposto é: necessidades de organização justificam funções de planejamento; condições de acesso justificam continuidade local; requisitos de interação e proteção delimitam sua operação; testes e registros delimitam quais conclusões podem ser sustentadas \citep{education_wolters2021,martin2024,eng06,OS18}.
Esse encadeamento organiza o estudo e não estabelece uma cadeia causal comprovada entre uso do aplicativo e desempenho acadêmico, mantendo distinta a hipótese educacional da observação técnica \citep{education_araka2020}.

Como regra analítica proposta, cada resultado responderá a um requisito dentro do ambiente executado: registro recuperado sustenta persistência naquele cenário, interface aberta sem rede sustenta disponibilidade nas condições testadas e dado inválido rejeitado sustenta aquela regra de validação \citep{eng06,eng11}.
Nenhum desses resultados, isoladamente, sustenta autonomia, aprendizagem, acessibilidade universal ou inclusão digital, pois tais construtos exigem outras medidas, públicos e desenhos de investigação \citep{education_araka2020,bong2021,martin2024}.

\label{sec:referencial-fim}
\clearpage

\section{Considerações finais}A execução anterior, v0.1.1, concluiu no GitHub Actions análise estática, seis testes Flutter, quatro testes de versionamento e compilação web. Commit: \texttt{ed3d39a}; execução: \texttt{37249382657}. A reabertura sem rede, a instalação PWA e a falha de quota continuam pendentes de teste específico. Esses resultados técnicos não demonstram eficácia pedagógica \citep{projeto}.

Na demonstração publicada da versão 0.1.0, a inspeção no navegador confirmou criação e conclusão de uma tarefa sintética e recuperação de seu título e estado após recarga online. Uma importação com versão 99 foi rejeitada com mensagem explícita; uma amostra válida restaurou duas tarefas, preservando os estados pendente e concluído. O indicador informou preparação do cache. Isso não comprova recarga sem conexão. A exportação exibiu solicitação de download, mas a captura do arquivo não foi confirmada nesta sessão. Os passos e limites constam de study/browser-observations.md. \citep{projeto}.

O artefato constitui pesquisa técnica de MVP, sem revisão por pares, avaliação pedagógica com estudantes, aprovação de banca ou certificação institucional. Não se inferem ganho de aprendizagem, usabilidade populacional ou disponibilidade universal dos testes técnicos. Faltam avaliação em dispositivos variados, sessões privadas, esgotamento de cota e estratégias robustas de atualização. Uma primeira visita sem conexão não deve ser assumida como suportada. A remoção dos dados do navegador pode eliminar tarefas; backup exportado é uma medida de recuperação, cuja integridade também precisa de teste, conforme os cenários ainda pendentes no registro técnico \citep{projeto}.
\subsection*{Declaração de autoria e uso de IA}Ferramentas de inteligência artificial auxiliaram a elaboração de código e texto nesta demonstração. A atribuição ``Estudantes — identificação omitida nesta demonstração'' não comprova autoria individual. Os estudantes devem conferir fontes, resultados e código, registrar suas contribuições e assumir a versão entregue. Aloisio é indicado como professor/orientador; sua revisão está pendente e não se presume autoria, assinatura ou aprovação; as responsabilidades e pendências estão registradas no projeto \citep{projeto}.
\subsection*{Limite de formatação}O modelo escolar fornecido exige Arial ou Times New Roman, tamanho 12, entrelinhas 1,5 e texto justificado. Nesta compilação XeLaTeX foi usada Liberation Sans por indisponibilidade de Arial. Essa substituição é declarada e deve ser validada pela escola; não se afirma conformidade tipográfica integral. As referências seguem apresentação autor/título/URL/data de acesso, sem auditoria completa de todas as normas ABNT; a política e suas pendências integram o registro do projeto \citep{projeto}.
\renewcommand{\refname}{4 Referências}\addcontentsline{toc}{section}{\protect\numberline{4}Referências}
\clearpage\section{Fundamentação rastreável e limites de transferência}

Esta síntese exploratória relaciona planejamento, interfaces, funcionamento local, engenharia e integridade da escrita. Foram consultados principalmente resumos e metadados; não foi realizada avaliação sistemática do risco de viés de cada estudo. As fontes sustentam os enunciados delimitados abaixo, sem validar o EstudaOffline nem comprovar aprendizagem. O registro de busca e seleção integra o protocolo do projeto \citep{projeto}.

\subsection{Planejamento e autorregulação da aprendizagem}

Revisão identificou uso continuado de instrumentos de sala presencial para medir autorregulação em ambientes digitais e lacunas em análises de registros. Período 2008–2018; medir atividades digitais não equivale a medir aprendizagem. \citep{education_araka2020}.


Revisão teórica discute suporte móvel oportuno a planejamento, monitoramento, reflexão e estratégias de aprendizagem. Não é ensaio de um planejador; oportunidades conceituais não são eficácia demonstrada. \citep{education_baars2022}.


Levantamento com 142 universitários encontrou associação negativa entre autorregulação e procrastinação acadêmica. Questionários e regressão não demonstram causalidade nem validade local para alunos do ADS. \citep{education_brahma2023}.


Em 73 universitários, condição combinando treinamento online e diários móveis apresentou maiores melhorias em componentes autorregulatórios. Intervenção combinada impede atribuir o resultado ao diário isolado ou a qualquer aplicativo. \citep{education_broadbent2020}.


Em duas turmas de 22 participantes, abordagem móvel gamificada melhorou desempenho, definição de metas e reflexão, mas não monitoramento. Amostra pequena; gamificação e instrução diferem de organizador local de tarefas. \citep{education_chen2024}.


Em 8.907 universitários durante a pandemia, motivação e procrastinação relacionaram-se à gestão do tempo e regulação do esforço. Contexto de emergência e relações observacionais não generalizam automaticamente para ensino técnico. \citep{education_forstervold2022}.


Em 1.216 universitários, autoeficácia autorregulatória e procrastinação apresentaram relações recíprocas; gestão do tempo foi dificuldade recorrente. Cursos assíncronos durante COVID-19; não testa aplicativo offline. \citep{education_han2023}.


Planejamento por smartphone e agenda associou-se a habilidades autorregulatórias; modelos apresentaram associações diferentes com notas. Associação negativa de smartphone e positiva de agenda sob controle estatístico não é causalidade nem avalia este aplicativo. \citep{education_hartley2022}.


Duas pesquisas-ação descrevem regulação não linear e construção compartilhada de conhecimento em aprendizagem móvel colaborativa. Contexto colaborativo qualitativo não valida planejador individual nem produz estimativa causal. \citep{education_lim2019}.


Em 450 universitários, estratégias metacognitivas e gestão do tempo foram examinadas como preditores de procrastinação. Desenho transversal não estabelece causalidade; contexto universitário específico. \citep{education_limone2020}.


Revisão de 38 estudos descreveu relação complexa entre aprendizagem móvel e autorregulação e recomendou integração curricular intencional. Síntese heterogênea de 2007–2019 não demonstra efeito de todo app móvel. \citep{education_palalas2020}.


Revisão de 107 estudos identificou planejamento, metas, priorização e organização de tarefas entre estratégias relevantes. Inclui educação e trabalho, artigos e dissertações; não demonstra efeito do MVP nem padroniza medida universal. \citep{education_patzak2025}.


Síntese de 31 revisões e meta-análises recomenda suportes autorregulatórios integrados ao longo do ciclo de aprendizagem. Revisões com contextos diversos; recomendação de desenho não constitui validação deste MVP. \citep{education_prasse2024}.


Em 272 universitários iranianos de inglês, maior autorregulação e controle atencional associaram-se a menor procrastinação. Amostra por conveniência e autorrelato; registro Consensus omitiu terceiro autor, corrigido pelo editor; 2024 online e 2026 fascículo. \citep{education_taghavi2024}.


Em 239 universitários de inglês como língua estrangeira, gestão do tempo e regulação do esforço associaram-se negativamente à procrastinação. Modelo de associações não demonstra eficácia de um planejador; população específica. \citep{education_tao2025}.


Em levantamento com 446 universitários, gestão do tempo contribuiu à predição de procrastinação além de variáveis motivacionais e estratégicas. Autorrelato e regressão não estabelecem efeito causal de um aplicativo. \citep{education_wolters2017}.


Revisão situa a gestão do tempo nas fases de antecipação, execução e avaliação da aprendizagem autorregulada. Síntese conceitual não comprova eficácia de aplicativo específico. \citep{education_wolters2021}.


Estudo identificou perfis ativo, desengajado e utilitário em app de prática de recuperação; autoeficácia e gestão do tempo predisseram notas. Perfis e associações não são efeito causal de uso de aplicativo; app inclui prática de recuperação ausente neste MVP. \citep{education_wong2026}.


Meta-análise encontrou efeito positivo moderado de intervenções de autorregulação no desempenho em ambientes online e híbridos. Intervenções e níveis escolares heterogêneos; efeito agrupado não valida EstudaOffline; 2022 online e 2023 fascículo. \citep{education_xu2022}.


Experimento com 60 graduandos encontrou melhores resultados e autorregulação em leitura de inglês apoiada por mecanismo móvel autorregulatório. Intervenção específica de leitura; 2016 online e 2018 fascículo; não transfere efeito ao EstudaOffline. \citep{education_zheng2016}.


\subsection{Acessibilidade, usabilidade e desigualdade digital}

Cobertura de rede, dispositivo e condições socioeconômicas afetam participação em ensino síncrono. Netnografia de tweets, com seleção de usuários de rede social; não representa todos os estudantes. \citep{azionya2021}.


Avaliação por pessoas cegas e surdas pode revelar diferenças na facilidade de uso de uma plataforma adaptada. Protótipo e população específicos; percepções não demonstram eficácia pedagógica ou acessibilidade universal. \citep{batanero2021}.


A revisão examina a formação docente para produção de materiais e ambientes digitais inclusivos; adotar apoio institucional no EstudaOffline é uma proposta de projeto. Revisão de 16 estudos; predominam entrevistas e surveys, com poucos indicadores objetivos. Não avalia EstudaOffline. \citep{bong2021}.


Acessibilidade, usabilidade e desenho universal para aprendizagem devem ser considerados em conjunto. Artigo conceitual com exemplos; não estima efeito causal nem certifica acessibilidade de um aplicativo. \citep{choi2024}.


Há lacunas na avaliação de efetividade de e-learning acessível a pessoas com dificuldades cognitivas. Revisão de estudos heterogêneos; resumo do editor consultado, texto integral não inspecionado. \citep{cinquin2019}.


A revisão organiza critérios de interface e sua relação com carga cognitiva; considerar tamanho de tela e desenho da interface no EstudaOffline é uma inferência de projeto. Revisão limitada a 2020–2022; não prova redução de carga cognitiva por uma interface específica. \citep{cob2023}.


O estudo avaliou recursos de voz e língua de sinais com estudantes e docentes; propor avaliação participativa do EstudaOffline é uma inferência de projeto. Estudo de métodos mistos de satisfação; não transfere resultados ao EstudaOffline nem à educação técnica brasileira. \citep{farhan2020}.


A avaliação de acessibilidade educacional deve considerar barreiras atitudinais além de aspectos técnicos. A síntese não permite afirmar que melhorias técnicas removem discriminação ou todas as barreiras. \citep{gobbo2023}.


Qualidade da conexão e condições de acesso integram a divisão digital no ensino online. Survey de estudantes chineses no contexto pandêmico; não autoriza inferência causal nem generalização direta ao Piauí. \citep{guo2022}.


Acesso à internet e habilidades digitais podem limitar continuidade educacional em universidades. Entrevistas por email podem excluir as pessoas com pior acesso; contexto paquistanês e pandêmico. \citep{jamil2023}.


Usabilidade de aplicações de aprendizagem móvel exige avaliação explícita e não pode ser presumida pela tecnologia. 23 publicações relevantes na revisão; publicado online em 2017, edição de 2018; não mede EstudaOffline. \citep{kumar2017}.


Diretrizes de usabilidade móvel devem ser atualizadas a partir de problemas observados. Síntese de 17 estudos e validação das diretrizes; condições e efeitos não podem ser transportados automaticamente. \citep{kumarguideline2019}.


Testes com usuários, observação e questionários são utilizados na avaliação de aprendizagem móvel. Mapeamento de 51 trabalhos de uma base; muitos testes controlados, limitando transferência para uso real. \citep{kumarmapping2024}.


Entrevistas com docentes e administradores descrevem estratégias de enfrentamento da divisão digital; combinar conectividade acessível, suporte e formação no projeto é uma proposta informada por esses relatos. 35 entrevistas em cinco universidades de Gana; relato de estratégias não é ensaio de eficácia causal. \citep{kumiyeboah2023}.


Existe pesquisa anterior sobre AVA com operação offline e uso em baixa conectividade, tornando inadequada uma alegação genérica de novidade de aplicativo offline educacional. Educa Offline é um AVA distinto de um planejador; o resumo e metadados não sustentam equivalência arquitetural, comparação de desempenho ou eficácia do EstudaOffline. SciSpace omitiu terceiro autor, corrigido com fonte primária. \citep{linhalis2026}.


Inclusão digital educacional depende de acesso, uso, fatores sociais e suporte organizacional. Revisão de 77 estudos em duas décadas; heterogeneidade e publicação online/volume posterior impedem uma estimativa única. \citep{martin2024}.


Acesso, práticas culturais e agência devem compor a análise de aprendizagem com conectividade limitada. Estudo em cinco países e no lockdown; gênero e responsabilidades domésticas dependem do contexto. \citep{mathrani2021}.


Facilidade de uso e acessibilidade integram fatores considerados na adoção de aprendizagem móvel. Revisão de 161 artigos, majoritariamente quantitativos; adoção e aprendizagem não são equivalentes. \citep{naveed2023}.


O artigo revisa padrões de interfaces móveis; aplicar sua discussão ao planejamento escolar do EstudaOffline é uma inferência de projeto. Revisão de interfaces móveis em geral; sua aplicação a planejamento escolar é uma inferência, não resultado direto. \citep{punchoojit2017}.


Prontidão digital dos estudantes e das escolas são dimensões distintas da desigualdade educacional. Dados e contextos internacionais; não demonstra que uma PWA resolve desigualdade estrutural. \citep{werfhorst2022}.


\subsection{Persistência local, cache, segurança e privacidade}

O estudo mostra que XSS na mesma origem pode manipular Cache API e ampliar o impacto sobre conteúdo armazenado. Cenários de ataque específicos; não demonstra vulnerabilidade no EstudaOffline nem estado atual de todos os navegadores. \citep{OS01}.


A sistematização reproduz ataques ligados a service workers e analisa medidas de mitigação dos navegadores. Panorama de 2022; mitigação atual exige consulta e teste específico. \citep{OS02}.


Os autores demonstram ataques de inferência de histórico explorando isolamento de service workers e cache. Resultados relativos às versões investigadas; não presumir que os ataques persistem em navegadores atuais. \citep{OS03}.


Experimentos mostram limitações dos controles e salvaguardas de cache nos navegadores examinados. Cache HTTP não equivale à Cache API; resultados históricos não certificam navegadores atuais. \citep{OS06}.


LoRe combina programação reativa e verificação para identificar interações que violam invariantes e coordená-las quando necessário. Garantias dependem do modelo, das propriedades especificadas e do compilador; não se transferem ao Flutter localStorage. \citep{OS07}.


Forking histories preserva restrições por histórico e expõe conflitos para reconciliação pelo usuário. Proposta de sincronização; não é resultado de adoção educacional nem regra necessária para app sem colaboração. \citep{OS08}.


O trabalho compara opções de criptografia para Web Storage e relata alternativas com impacto de desempenho limitado no cenário avaliado. Criptografia local não elimina XSS nem valida gestão de chaves; dados e implementação do MVP não foram avaliados. \citep{OS09}.


BeauForT propõe consenso e sincronização no navegador para tolerar réplicas bizantinas e redes instáveis. Sistema descentralizado com múltiplas partes; não constitui proteção pronta para aplicativo local sem sincronização. \citep{OS10}.


A análise de 137 apps infantis e entrevistas com 12 desenvolvedores identificaram barreiras ao uso de SDKs favoráveis à privacidade. Estudo qualitativo e ferramenta de apoio; não demonstra que remover SDKs garante segurança integral. \citep{OS13}.


A análise de 42 apps Android de cuidado infantil e seus backends identificou rastreadores e riscos de exposição de dados. Contexto de creches e backends móveis; não extrapolar vulnerabilidades para app local com dados sintéticos. \citep{OS14}.


Oficinas de codesign orientaram KOALA Hero para tornar riscos de tratamento de dados compreensíveis a crianças. Os autores deixam avaliação empírica futura em aberto; não evidencia mudança comprovada de comportamento. \citep{OS15}.


Análise estática e dinâmica de 20.195 apps identificou práticas de rastreamento e permissões que divergiam de políticas de proteção infantil. Amostra Google Play e políticas do período; não mede riscos de uma PWA específica nem todos os ecossistemas. \citep{OS17}.


A revisão selecionou 21 de 64 trabalhos e apontou lacunas na literatura sobre segurança de apps educacionais. Revisão secundária; apenas resumo consultado, sem auditoria de busca ou qualidade de cada estudo. \citep{OS18}.


\subsection{Engenharia, testes e reprodutibilidade}

A revisão de 103 estudos distingue testes funcionais e não funcionais e identifica desafios do ecossistema Android. Foco Android e literatura até 2016; não valida Flutter web nem este aplicativo. \citep{eng01}.


A revisão de 56 artigos identifica necessidades de adaptação dos frameworks de automação ao domínio móvel. Arquitetura MATF proposta, com implementação futura; não comprova eficácia do MVP. \citep{eng02}.


O mapeamento de 131 artigos organiza técnicas, atividades de teste e formas de avaliação, com predominância de Android. Estudo secundário; o ano online 2018 difere do volume de 2019. Não transferir resultados Android diretamente à web. \citep{eng03}.


O estudo terciário organiza 24 revisões sobre engenharia de aplicativos e identifica lacunas em manutenção e plataformas cruzadas. Síntese de revisões; não demonstra qualidade ou efeito pedagógico do aplicativo. \citep{eng04}.


O mapeamento de 79 estudos identifica necessidade de elicitar requisitos de teste cedo e avaliar em contextos reais. Mapeamento secundário de literatura anterior; não substitui ensaios atuais no navegador. \citep{eng05}.


O mapeamento de 56 estudos caracteriza técnicas e ferramentas de teste dinâmico de requisitos não funcionais. Classifica trabalhos existentes; segurança e confiabilidade do MVP exigem testes específicos. \citep{eng06}.


O artigo propõe explicitar o artefato e o objetivo de pesquisa para tornar mais precisa a investigação em design science. Validação metodológica em artigos DSR; não é medida de eficácia educacional. \citep{eng07}.


A reavaliação de um método para estudos de repositórios mostra utilidade de caracterizar a reprodutibilidade além da disponibilidade. Escopo de mineração de repositórios; não garante reprodução do MVP. \citep{eng08}.


A revisão e reexecução de modelos de deep learning mostram que compartilhar código não basta para reproduzir resultados. Escopo deep learning; o aplicativo não usa IA. Relevância limitada a documentação e pacotes reprodutíveis. \citep{eng10}.


O mapeamento de 40 estudos identifica dependência do domínio nas ferramentas e práticas de reprodução experimental. Ferramentas dependem do domínio; não autoriza generalização direta para estudantes. \citep{eng11}.


Sete replicações externas identificam dificuldades relacionadas a instruções pouco claras e defeitos nos artefatos. Experimento em sistemas configuráveis; aplicar como orientação de documentação, não como resultado do MVP. \citep{eng12}.


A revisão de replicações identifica relato incompleto e problemas para interpretar resultados confirmatórios. Escopo previsão de esforço e pair programming; relato do resumo não estabelece validade de qualquer replicação. \citep{eng13}.


A revisão de 69 artigos organiza ferramentas e desafios de integração, entrega e implantação contínuas. Revisão até junho de 2016; resultados exigem contextualização em pipelines atuais. \citep{eng14}.


A revisão de 101 estudos empíricos identifica benefícios e desafios técnicos associados à integração contínua. Estudos heterogêneos; benefícios de CI não equivalem a ausência de bugs ou eficácia pedagógica. \citep{eng15}.


O estudo de três empresas mostra que a implementação de CI varia com o contexto e pode criar restrições ao feedback. Três empresas; generalização limitada e ausência de avaliação deste pipeline. \citep{eng16}.


O FEDS separa avaliações formativas e somativas, artificiais e naturalistas, e orienta o planejamento dos episódios. Framework metodológico com exemplos; não certifica o artefato nem substitui avaliação com usuários. \citep{eng17}.


A proposta de estrutura de artefatos para experimentos humanos organiza elementos alinhados às políticas de badges ACM. Exige validação externa adicional; um repositório estruturado não recebe automaticamente badge ACM. \citep{eng18}.


O MEDS propõe planejar objetivos, estratégias, propriedades e episódios de avaliação para ajustar rigor e recursos em DSR. Método orientador; não demonstra qualidade do MVP nem efeitos nos estudantes. \citep{eng20}.


\subsection{Integridade, atribuição e revisão da escrita}

PRISMA 2020 fornece checklist de 27 itens e fluxogramas para relatar revisões sistemáticas de modo transparente. Diretriz de relato, não valida qualidade metodológica nem transforma esta busca exploratória em revisão sistemática. \citep{SU01}.


O artigo de explicação detalha a justificativa e exemplifica recomendações de cada item PRISMA 2020. Publicação complementar à mesma diretriz: fonte distinta, não estudo independente de eficácia. \citep{SU02}.


PRISMA-S apresenta 16 itens para descrever fontes e procedimentos de busca de forma reproduzível. Checklist desenvolvido por consenso; completude do relato não garante exaustividade ou ausência de viés. \citep{SU03}.


O manifesto propõe melhorar métodos, transparência, reprodução, avaliação e incentivos na pesquisa. Proposta baseada em evidências variadas; adoção ampla ainda requer avaliação iterativa e não certifica o MVP. \citep{SU04}.


A revisão de 239 estudos identifica avanços na detecção de plágio e lacunas de avaliação rigorosa dos sistemas. Literatura de 2013–2018; detecção computacional não substitui julgamento contextual nem garante originalidade. \citep{SU05}.


A avaliação de 636 referências geradas por GPT-3.5 e GPT-4 encontrou referências fabricadas e erros em referências reais. Modelos e prompts de abril de 2023; taxas não são estimativas para modelos atuais nem para esta pesquisa com verificação externa. \citep{SU06}.


Nos 30 textos médicos gerados por GPT-3.5, os autores encontraram referências fabricadas ou com metadados incorretos. 115 referências e prompts médicos específicos; não generalizar taxas para todos os gêneros ou modelos. \citep{SU07}.


AcPgChecker compara documentos offline com similaridade de cosseno e um limiar para sinalizar possível plágio. Resumo não apresenta métricas completas; limiar de similaridade não constitui decisão ética válida por si só. \citep{SU08}.


A orientação metodológica organiza pôsteres em mensagem central, apresentação visual e comunicação oral ajustadas ao público, em vez de reproduzir todo o artigo no painel principal. Guia voltado à pesquisa em saúde, não experimento de validação de um template universal; usado como fundamento editorial anterior à janela de 2016–2026, sem inferir eficácia educacional. \citep{SU09}.

\clearpage\begin{thebibliography}{99}

\bibitem[Ahmad Faudzi et al.(2023)]{cob2023} AHMAD FAUDZI, M.; CHE COB, Zaihisma; OMAR, Ridha; SHARUDIN, Sharul Azim; GHAZALI, Masitah. Investigating the User Interface Design Frameworks of Current Mobile Learning Applications: A Systematic Review. \textbf{Education Sciences}, 2023. DOI: 10.3390/educsci13010094. URL: \url{https://doi.org/10.3390/educsci13010094}. Consulta: 5 out. 2026. Página do trecho não conferida; nível de acesso registrado no fichamento.

\bibitem[Anchundia e Fonseca C.(2020)]{eng11} ANCHUNDIA, Carlos E.; FONSECA C., Efraín R.. Resources for Reproducibility of Experiments in Empirical Software Engineering: Topics Derived From a Secondary Study. \textbf{IEEE Access}, 2020. DOI: 10.1109/ACCESS.2020.2964587. URL: \url{https://doi.org/10.1109/ACCESS.2020.2964587}. Consulta: 5 out. 2026. Página do trecho não conferida; nível de acesso registrado no fichamento.

\bibitem[Araka et al.(2020)]{education_araka2020} ARAKA, Eric; MAINA, E.; GITONGA, Rhoda; OBOKO, Robert O.. Research trends in measurement and intervention tools for self-regulated learning for e-learning environments—systematic review (2008–2018). \textbf{Research and Practice in Technology Enhanced Learning}, 2020. DOI: 10.1186/s41039-020-00129-5. URL: \url{https://doi.org/10.1186/s41039-020-00129-5}. Consulta: 5 out. 2026. Localizador utilizado: p. 4, 15–16. Leitura: \url{https://link.springer.com/content/pdf/10.1186/s41039-020-00129-5.pdf}.

\bibitem[Azionya e Nhedzi(2021)]{azionya2021} AZIONYA, C.; NHEDZI, Abyshey. THE DIGITAL DIVIDE AND HIGHER EDUCATION CHALLENGE WITH EMERGENCY ONLINE LEARNING: ANALYSIS OF TWEETS IN THE WAKE OF THE COVID-19 LOCKDOWN. \textbf{Turkish Online Journal of Distance Education}, 2021. DOI: 10.17718/tojde.1002822. URL: \url{https://doi.org/10.17718/tojde.1002822}. Consulta: 5 out. 2026. Página do trecho não conferida; nível de acesso registrado no fichamento.

\bibitem[Baars e Viberg(2022)]{education_baars2022} BAARS, Martine; VIBERG, Olga. Mobile Learning to Support Self-Regulated Learning: A Theoretical Review. \textbf{Int. J. Mob. Blended Learn.}, 2022. DOI: 10.4018/IJMBL.315628. URL: \url{https://doi.org/10.4018/IJMBL.315628}. Consulta: 5 out. 2026. Página do trecho não conferida; nível de acesso registrado no fichamento.

\bibitem[Baskerville et al.(2026)]{eng20} BASKERVILLE, Richard; PRIES-HEJE, Jan; VENABLE, John R.. MEDS: Methodology for Evaluation in Design Science. \textbf{European Journal of Information Systems}, 2026. DOI: 10.1080/0960085X.2026.2627280. URL: \url{https://doi.org/10.1080/0960085X.2026.2627280}. Consulta: 5 out. 2026. Página do trecho não conferida; nível de acesso registrado no fichamento.

\bibitem[Batanero-Ochaíta et al.(2021)]{batanero2021} BATANERO-OCHAÍTA, Concepción; DE-MARCOS, Luis; RIVERA, Luis Felipe; HOLVIKIVI, J.; HILERA, J.; TORTOSA, Salvador Otón. Improving Accessibility in Online Education: Comparative Analysis of Attitudes of Blind and Deaf Students Toward an Adapted Learning Platform. \textbf{IEEE Access}, 2021. DOI: 10.1109/ACCESS.2021.3095041. URL: \url{https://doi.org/10.1109/ACCESS.2021.3095041}. Consulta: 5 out. 2026. Página do trecho não conferida; nível de acesso registrado no fichamento.

\bibitem[Berihun et al.(2023)]{eng02} BERIHUN, Natnael Gonfa; DONGMO, Cyrille; POLL, J. A. van der. The Applicability of Automated Testing Frameworks for Mobile Application Testing: A Systematic Literature Review. \textbf{Computers}, 2023. DOI: 10.3390/computers12050097. URL: \url{https://doi.org/10.3390/computers12050097}. Consulta: 5 out. 2026. Página do trecho não conferida; nível de acesso registrado no fichamento.

\bibitem[Bhattacharyya et al.(2023)]{SU07} BHATTACHARYYA, Mehul; MILLER, Valerie M.; BHATTACHARYYA, Debjani; MILLER, Larry E.. High Rates of Fabricated and Inaccurate References in ChatGPT-Generated Medical Content. \textbf{Cureus}, 2023. DOI: 10.7759/cureus.39238. URL: \url{https://doi.org/10.7759/cureus.39238}. Consulta: 5 out. 2026. Página do trecho não conferida; nível de acesso registrado no fichamento.

\bibitem[Bong e Chen(2021)]{bong2021} BONG, W. K.; CHEN, Wei-Qin. Increasing faculty’s competence in digital accessibility for inclusive education: a systematic literature review. \textbf{International Journal of Inclusive Education}, 2021. DOI: 10.1080/13603116.2021.1937344. URL: \url{https://doi.org/10.1080/13603116.2021.1937344}. Consulta: 5 out. 2026. Página do trecho não conferida; nível de acesso registrado no fichamento.

\bibitem[Brahma e Saikia(2023)]{education_brahma2023} BRAHMA, B.; SAIKIA, P.. Influence of self-regulated learning on the academic procrastination of college students. \textbf{Journal of Education and Health Promotion}, 2023. DOI: 10.4103/jehp.jehp\_1106\_22. URL: \url{https://doi.org/10.4103/jehp.jehp_1106_22}. Consulta: 5 out. 2026. Página do trecho não conferida; nível de acesso registrado no fichamento.

\bibitem[Broadbent et al.(2020)]{education_broadbent2020} BROADBENT, Jaclyn; PANADERO, E.; FULLER-TYSZKIEWICZ, M.. Effects of mobile-app learning diaries vs online training on specific self-regulated learning components. \textbf{Educational Technology Research and Development}, 2020. DOI: 10.1007/s11423-020-09781-6. URL: \url{https://doi.org/10.1007/s11423-020-09781-6}. Consulta: 5 out. 2026. Página do trecho não conferida; nível de acesso registrado no fichamento.

\bibitem[Chen et al.(2024)]{education_chen2024} CHEN, Yi-Chun; HWANG, Gwo-Jen; LAI, Chiu-Lin. Motivating students to become self-regulatory learners: A gamified mobile self-regulated learning approach. \textbf{Education and Information Technologies}, 2024. DOI: 10.1007/s10639-024-12462-z. URL: \url{https://doi.org/10.1007/s10639-024-12462-z}. Consulta: 5 out. 2026. Página do trecho não conferida; nível de acesso registrado no fichamento.

\bibitem[Choi e Seo(2024)]{choi2024} CHOI, Gi Woong; SEO, JooYoung. Accessibility, Usability, and Universal Design for Learning: Discussion of Three Key LX/UX Elements for Inclusive Learning Design. \textbf{TechTrends}, 2024. DOI: 10.1007/s11528-024-00987-6. URL: \url{https://doi.org/10.1007/s11528-024-00987-6}. Consulta: 5 out. 2026. Página do trecho não conferida; nível de acesso registrado no fichamento.

\bibitem[Cinquin et al.(2019)]{cinquin2019} CINQUIN, Pierre-Antoine; GUITTON, P.; SAUZÉON, H.. Online e-learning and cognitive disabilities: A systematic review. \textbf{Comput. Educ.}, 2019. DOI: 10.1016/j.compedu.2018.12.004. URL: \url{https://doi.org/10.1016/j.compedu.2018.12.004}. Consulta: 5 out. 2026. Página do trecho não conferida; nível de acesso registrado no fichamento.

\bibitem[Dehling et al.(2019)]{OS06} DEHLING, Florian; MENGEL, Tobias; IACONO, Luigi Lo. Rotten Cellar: Security and Privacy of the Browser Cache Revisited. \textbf{Secure IT Systems: 24th Nordic Conference, NordSec 2019, pp.20–36}, 2019. DOI: 10.1007/978-3-030-35055-0\_2. URL: \url{https://doi.org/10.1007/978-3-030-35055-0_2}. Consulta: 5 out. 2026. Página do trecho não conferida; nível de acesso registrado no fichamento.

\bibitem[Dipongkor et al.(2021)]{SU08} DIPONGKOR, Atish Kumar; ISLAM, Rayhanul; SHAFIUZZAMAN, Md.; NASHIRY, Asif; GALIB, Syed Md.; MAZUMDER, Khaza Moinuddin. AcPgChecker: Detection of Plagiarism among Academic and Scientific Writings. \textbf{2021 Joint 10th International Conference on Informatics, Electronics \& Vision (ICIEV) and 2021 5th International Conference on Imaging, Vision \& Pattern Recognition (icIVPR)}, 2021. DOI: 10.1109/icievicivpr52578.2021.9564123. URL: \url{https://doi.org/10.1109/icievicivpr52578.2021.9564123}. Consulta: 5 out. 2026. Página do trecho não conferida; nível de acesso registrado no fichamento.

\bibitem[Ekambaranathan et al.(2023)]{OS13} EKAMBARANATHAN, Anirudh; ZHAO, Jun; CHALHOUB, George. Navigating the Data Avalanche: Towards Supporting Developers in Developing Privacy-Friendly Children’s Apps. \textbf{Proceedings of the ACM on interactive, mobile, wearable and ubiquitous technologies}, 2023. DOI: 10.1145/3596267. URL: \url{https://doi.org/10.1145/3596267}. Consulta: 5 out. 2026. Página do trecho não conferida; nível de acesso registrado no fichamento.

\bibitem[Elazhary et al.(2021)]{eng16} ELAZHARY, Omar; WERNER, Colin M.; LI, Ze-Shi; LOWLIND, Derek; ERNST, Neil A.; STOREY, M.. Uncovering the Benefits and Challenges of Continuous Integration Practices. \textbf{IEEE Transactions on Software Engineering}, 2021. DOI: 10.1109/TSE.2021.3064953. URL: \url{https://doi.org/10.1109/TSE.2021.3064953}. Consulta: 5 out. 2026. Página do trecho não conferida; nível de acesso registrado no fichamento.

\bibitem[Farhan e Razmak(2020)]{farhan2020} FARHAN, Wejdan; RAZMAK, Jamil. A comparative study of an assistive e-learning interface among students with and without visual and hearing impairments. \textbf{Disability and Rehabilitation: Assistive Technology}, 2020. DOI: 10.1080/17483107.2020.1786733. URL: \url{https://doi.org/10.1080/17483107.2020.1786733}. Consulta: 5 out. 2026. Página do trecho não conferida; nível de acesso registrado no fichamento.

\bibitem[Flutter(2026)]{flutter} Flutter. Web FAQ. \url{https://docs.flutter.dev/platform-integration/web/faq}. Acesso em 5 out. 2026. Registro técnico, fora da contagem de fontes científicas.

\bibitem[Foltýnek et al.(2019)]{SU05} FOLTÝNEK, Tomáš; MEUSCHKE, Norman; GIPP, Bela. Academic Plagiarism Detection: A Systematic Literature Review. \textbf{ACM Computing Surveys}, 2019. DOI: 10.1145/3345317. URL: \url{https://doi.org/10.1145/3345317}. Consulta: 5 out. 2026. Página do trecho não conferida; nível de acesso registrado no fichamento.

\bibitem[Forstervold et al.(2022)]{education_forstervold2022} FORSTERVOLD, Knut Inge; LUDVIGSEN, S.; STRØMSØ, Helge I.. Students’ time management and procrastination in the wake of the pandemic. \textbf{Educational Psychology}, 2022. DOI: 10.1080/01443410.2022.2102582. URL: \url{https://doi.org/10.1080/01443410.2022.2102582}. Consulta: 5 out. 2026. Página do trecho não conferida; nível de acesso registrado no fichamento.

\bibitem[Gobbo e Bezerra(2023)]{gobbo2023} GOBBO, José Alcides; BEZERRA, B.. Emerging Themes for Digital Accessibility in Education. \textbf{Sustainability}, 2023. DOI: 10.3390/su151411392. URL: \url{https://doi.org/10.3390/su151411392}. Consulta: 5 out. 2026. Página do trecho não conferida; nível de acesso registrado no fichamento.

\bibitem[Gonzalez-Barahona e Robles(2023)]{eng08} GONZALEZ-BARAHONA, Jesus M.; ROBLES, G.. Revisiting the reproducibility of empirical software engineering studies based on data retrieved from development repositories. \textbf{Inf. Softw. Technol.}, 2023. DOI: 10.1016/j.infsof.2023.107318. URL: \url{https://doi.org/10.1016/j.infsof.2023.107318}. Consulta: 5 out. 2026. Página do trecho não conferida; nível de acesso registrado no fichamento.

\bibitem[Gruber et al.(2022)]{OS14} GRUBER, Moritz; HÖFIG, Christian; GOLLA, Maximilian; URBAN, Tobias; GROSSE-KAMPMANN, Matteo. "We may share the number of diaper changes": A Privacy and Security Analysis of Mobile Child Care Applications. \textbf{Proceedings on Privacy Enhancing Technologies}, 2022. DOI: 10.56553/popets-2022-0078. URL: \url{https://doi.org/10.56553/popets-2022-0078}. Consulta: 5 out. 2026. Página do trecho não conferida; nível de acesso registrado no fichamento.

\bibitem[Guevara-Vega et al.(2024)]{eng18} GUEVARA-VEGA, C.; BERNÁRDEZ, B.; CRUZ, Margarita; DURÁN, A.; RUIZ-CORTÉS, Antonio; SOLARI, Martin. Research artifacts for human-oriented experiments in software engineering: An ACM badges-driven structure proposal. \textbf{J. Syst. Softw.}, 2024. DOI: 10.1016/j.jss.2024.112187. URL: \url{https://doi.org/10.1016/j.jss.2024.112187}. Consulta: 5 out. 2026. Página do trecho não conferida; nível de acesso registrado no fichamento.

\bibitem[Guo e Wan(2022)]{guo2022} GUO, Congbin; WAN, Boshen. The digital divide in online learning in China during the COVID-19 pandemic. \textbf{Technology in Society}, 2022. DOI: 10.1016/j.techsoc.2022.102122. URL: \url{https://doi.org/10.1016/j.techsoc.2022.102122}. Consulta: 5 out. 2026. Página do trecho não conferida; nível de acesso registrado no fichamento.

\bibitem[Haas et al.(2023)]{OS07} HAAS, Julian; MOGK, Ragnar; YANAKIEVA, Elena; BIENIUSA, Annette; MEZINI, Mira. LoRe: A Programming Model for Verifiably Safe Local-First Software. \textbf{ACM Transactions on Programming Languages and Systems}, 2023. DOI: 10.1145/3633769. URL: \url{https://doi.org/10.1145/3633769}. Consulta: 5 out. 2026. Página do trecho não conferida; nível de acesso registrado no fichamento.

\bibitem[Han et al.(2023)]{education_han2023} HAN, Jaeyun; DIGIACOMO, D.; USHER, Ellen L.. College students’ self-regulation in asynchronous online courses during COVID-19. \textbf{Studies in Higher Education}, 2023. DOI: 10.1080/03075079.2023.2201608. URL: \url{https://doi.org/10.1080/03075079.2023.2201608}. Consulta: 5 out. 2026. Página do trecho não conferida; nível de acesso registrado no fichamento.

\bibitem[Hartley et al.(2022)]{education_hartley2022} HARTLEY, Kendall; SHREVE, Emily; GIANOUTSOS, Dan; BENDIXEN, Lisa D.. Smartphone as a Self-regulatory Planning Tool: Promise or Peril. \textbf{International journal of interactive mobile technologies}, 2022. DOI: 10.3991/ijim.v16i14.28783. URL: \url{https://doi.org/10.3991/ijim.v16i14.28783}. Consulta: 5 out. 2026. Página do trecho não conferida; nível de acesso registrado no fichamento.

\bibitem[Jamil e Muschert(2023)]{jamil2023} JAMIL, S.; MUSCHERT, G.. The COVID-19 Pandemic and E-Learning: The Digital Divide and Educational Crises in Pakistan’s Universities. \textbf{The American Behavioral Scientist}, 2023. DOI: 10.1177/00027642231156779. URL: \url{https://doi.org/10.1177/00027642231156779}. Consulta: 5 out. 2026. Página do trecho não conferida; nível de acesso registrado no fichamento.

\bibitem[Jannes et al.(2023)]{OS10} JANNES, Kristof; BENI, Emad Heydari; LAGAISSE, Bert; JOOSEN, Wouter. BeauForT: Robust Byzantine Fault Tolerance for Client-Centric Mobile Web Applications. \textbf{IEEE Transactions on Parallel and Distributed Systems}, 2023. DOI: 10.1109/TPDS.2023.3241963. URL: \url{https://doi.org/10.1109/TPDS.2023.3241963}. Consulta: 5 out. 2026. Página do trecho não conferida; nível de acesso registrado no fichamento.

\bibitem[Júnior et al.(2022)]{eng06} JÚNIOR, Misael C.; AMALFITANO, Domenico; GARCÉS, Lina; FASOLINO, A.; ANDRADE, Stevão A.; DELAMARO, M.. Dynamic Testing Techniques of Non-functional Requirements in Mobile Apps: A Systematic Mapping Study. \textbf{ACM Computing Surveys (CSUR)}, 2022. DOI: 10.1145/3507903. URL: \url{https://doi.org/10.1145/3507903}. Consulta: 5 out. 2026. Página do trecho não conferida; nível de acesso registrado no fichamento.

\bibitem[Karami et al.(2021)]{OS03} KARAMI, Soroush; ILIA, Panagiotis; POLAKIS, Jason. Awakening the Web's Sleeper Agents: Misusing Service Workers for Privacy Leakage.. \textbf{Network and Distributed System Security Symposium (NDSS), 2021}, 2021. DOI: 10.14722/NDSS.2021.23104. URL: \url{https://doi.org/10.14722/NDSS.2021.23104}. Consulta: 5 out. 2026. Página do trecho não conferida; nível de acesso registrado no fichamento.

\bibitem[Kong et al.(2019)]{eng01} KONG, Pingfan; LI, Li; GAO, Jun; LIU, Kui; BISSYANDÉ, Tegawendé F.; KLEIN, Jacques. Automated Testing of Android Apps: A Systematic Literature Review. \textbf{IEEE Transactions on Reliability}, 2019. DOI: 10.1109/TR.2018.2865733. URL: \url{https://doi.org/10.1109/TR.2018.2865733}. Consulta: 5 out. 2026. Página do trecho não conferida; nível de acesso registrado no fichamento.

\bibitem[Kumar e Mohite(2017)]{kumar2017} KUMAR, B.; MOHITE, P.. Usability of mobile learning applications: a systematic literature review. \textbf{Journal of Computers in Education}, 2017. DOI: 10.1007/s40692-017-0093-6. URL: \url{https://doi.org/10.1007/s40692-017-0093-6}. Consulta: 5 out. 2026. Página do trecho não conferida; nível de acesso registrado no fichamento.

\bibitem[Kumar et al.(2019)]{kumarguideline2019} KUMAR, B.; GOUNDAR, M. S.; CHAND, S.. Usability guideline for Mobile learning applications: an update. \textbf{Education and Information Technologies}, 2019. DOI: 10.1007/s10639-019-09937-9. URL: \url{https://doi.org/10.1007/s10639-019-09937-9}. Consulta: 5 out. 2026. Página do trecho não conferida; nível de acesso registrado no fichamento.

\bibitem[Kumar et al.(2024)]{kumarmapping2024} KUMAR, B.; CHAND, S.; GOUNDAR, M. S.. Usability testing of mobile learning applications: a systematic mapping study. \textbf{The International Journal of Information and Learning Technology}, 2024. DOI: 10.1108/IJILT-03-2023-0029. URL: \url{https://doi.org/10.1108/IJILT-03-2023-0029}. Consulta: 5 out. 2026. Página do trecho não conferida; nível de acesso registrado no fichamento.

\bibitem[Kumi-Yeboah et al.(2023)]{kumiyeboah2023} KUMI-YEBOAH, Alex; KIM, YangHyun; ARMAH, Yaa Essah. Strategies for overcoming the digital divide during the COVID-19 pandemic in higher education institutions in Ghana. \textbf{Br. J. Educ. Technol.}, 2023. DOI: 10.1111/bjet.13356. URL: \url{https://doi.org/10.1111/bjet.13356}. Consulta: 5 out. 2026. Página do trecho não conferida; nível de acesso registrado no fichamento.

\bibitem[Lim et al.(2019)]{education_lim2019} LIM, Genevieve; SHELLEY, A.; HEO, Dongcheol. The regulation of learning and co-creation of new knowledge in mobile learning. \textbf{Knowledge Management \& E-Learning: An International Journal}, 2019. DOI: 10.34105/j.kmel.2019.11.024. URL: \url{https://doi.org/10.34105/j.kmel.2019.11.024}. Consulta: 5 out. 2026. Página do trecho não conferida; nível de acesso registrado no fichamento.

\bibitem[Limone et al.(2020)]{education_limone2020} LIMONE, P.; SINATRA, M.; CEGLIE, F.; MONACIS, L.. Examining Procrastination among University Students through the Lens of the Self-Regulated Learning Model. \textbf{Behavioral Sciences}, 2020. DOI: 10.3390/bs10120184. URL: \url{https://doi.org/10.3390/bs10120184}. Consulta: 5 out. 2026. Página do trecho não conferida; nível de acesso registrado no fichamento.

\bibitem[Linhalis et al.(2026)]{linhalis2026} LINHALIS, Flávia; PAVÃO, Thiago Máximo; SILVA, André Constantino da. Educa Offline: concepción de un Ambiente Virtual de Aprendizaje (AVA) en contextos de baja conectividad y vulnerabilidad social. \textbf{Revista Latinoamericana de Tecnología Educativa - RELATEC}, 2026. DOI: 10.17398/1695-288X.25.1.59. URL: \url{https://doi.org/10.17398/1695-288X.25.1.59}. Consulta: 5 out. 2026. Página do trecho não conferida; nível de acesso registrado no fichamento.

\bibitem[Litt e Jackson(2022)]{OS08} LITT, Geoffrey; JACKSON, B.. Merge what you can, fork what you can't: managing data integrity in local-first software. \textbf{9th Workshop on Principles and Practice of Consistency for Distributed Data (PaPoC), 2022}, 2022. DOI: 10.1145/3517209.3524041. URL: \url{https://doi.org/10.1145/3517209.3524041}. Consulta: 5 out. 2026. Página do trecho não conferida; nível de acesso registrado no fichamento.

\bibitem[Liu et al.(2022)]{eng10} LIU, Chao; GAO, Cui-Yun; XIA, Xin; LO, David; GRUNDY, John C.; YANG, Xiao-Hu. On the Reproducibility and Replicability of Deep Learning in Software Engineering. \textbf{ACM Trans. Softw. Eng. Methodol.}, 2022. DOI: 10.1145/3477535. URL: \url{https://doi.org/10.1145/3477535}. Consulta: 5 out. 2026. Página do trecho não conferida; nível de acesso registrado no fichamento.

\bibitem[MacFadden e Qiu(2022)]{OS09} MACFADDEN, Michael; QIU, Meikang. Performance Impacts of JavaScript-Based Encryption of HTML5 Web Storage for Enhanced Privacy. \textbf{IEEE SmartCloud 2022}, 2022. DOI: 10.1109/SmartCloud55982.2022.00037. URL: \url{https://doi.org/10.1109/SmartCloud55982.2022.00037}. Consulta: 5 out. 2026. Página do trecho não conferida; nível de acesso registrado no fichamento.

\bibitem[Martin et al.(2024)]{martin2024} MARTIN, Florence; CEVIKER, E.; GEZER, T.. From digital divide to digital equity: Systematic review of two decades of research on educational digital divide factors, dimensions, and interventions. \textbf{Journal of Research on Technology in Education}, 2024. DOI: 10.1080/15391523.2024.2425442. URL: \url{https://doi.org/10.1080/15391523.2024.2425442}. Consulta: 5 out. 2026. Página do trecho não conferida; nível de acesso registrado no fichamento.

\bibitem[Mathrani et al.(2021)]{mathrani2021} MATHRANI, A.; SARVESH, Tarushikha; UMER, R.. Digital divide framework: online learning in developing countries during the COVID-19 lockdown. \textbf{Globalisation, Societies and Education}, 2021. DOI: 10.1080/14767724.2021.1981253. URL: \url{https://doi.org/10.1080/14767724.2021.1981253}. Consulta: 5 out. 2026. Página do trecho não conferida; nível de acesso registrado no fichamento.

\bibitem[MDN(2026a)]{local} MDN. Window: localStorage property. \url{https://developer.mozilla.org/en-US/docs/Web/API/Window/localStorage}. Acesso em 5 out. 2026. Registro técnico, fora da contagem de fontes científicas.

\bibitem[MDN(2026b)]{quota} MDN. Storage quotas and eviction criteria. \url{https://developer.mozilla.org/en-US/docs/Web/API/Storage_API/Storage_quotas_and_eviction_criteria}. Acesso em 5 out. 2026. Registro técnico, fora da contagem de fontes científicas.

\bibitem[MDN(2026c)]{sw} MDN. Service Worker API. \url{https://developer.mozilla.org/en-US/docs/Web/API/Service_Worker_API}. Acesso em 5 out. 2026. Registro técnico, fora da contagem de fontes científicas.

\bibitem[Mkpojiogu et al.(2021)]{OS18} MKPOJIOGU, Emmanuel O.C.; HUSSAIN, Azham; AGBUDU, Monday Onah. Security Issues in the Use of Mobile Educational Apps: A Review. \textbf{International Journal of Interactive Mobile Technologies (ijim)}, 2021. DOI: 10.3991/IJIM.V15I06.20631. URL: \url{https://doi.org/10.3991/IJIM.V15I06.20631}. Consulta: 5 out. 2026. Página do trecho não conferida; nível de acesso registrado no fichamento.

\bibitem[Munafò et al.(2017)]{SU04} MUNAFÒ, Marcus R.; NOSEK, Brian A.; BISHOP, Dorothy V. M.; BUTTON, Katherine S.; CHAMBERS, Christopher D.; SERT, Nathalie Percie du; SIMONSOHN, Uri; WAGENMAKERS, Eric-Jan; WARE, Jennifer J.; IOANNIDIS, John P. A.. A manifesto for reproducible science. \textbf{Nature Human Behaviour}, 2017. DOI: 10.1038/s41562-016-0021. URL: \url{https://doi.org/10.1038/s41562-016-0021}. Consulta: 5 out. 2026. Página do trecho não conferida; nível de acesso registrado no fichamento.

\bibitem[Naveed et al.(2023)]{naveed2023} NAVEED, Q.; CHOUDHARY, Heena; AHMAD, Naim; ALQAHTANI, Jarallah; QAHMASH, A.. Mobile Learning in Higher Education: A Systematic Literature Review. \textbf{Sustainability}, 2023. DOI: 10.3390/su151813566. URL: \url{https://doi.org/10.3390/su151813566}. Consulta: 5 out. 2026. Página do trecho não conferida; nível de acesso registrado no fichamento.

\bibitem[Page et al.(2021a)]{SU01} PAGE, Matthew J.; MCKENZIE, Joanne E.; BOSSUYT, Patrick M.M.; BOUTRON, Isabelle; HOFFMANN, Tammy; MULROW, Cynthia D.; SHAMSEER, Larissa; TETZLAFF, Jennifer; AKL, Elie A.; BRENNAN, Sue E.; CHOU, Roger; GLANVILLE, Julie; GRIMSHAW, Jeremy M.; HRÓBJARTSSON, Asbjørn; LALU, Manoj M.; LI, Tianjing; LODER, Elizabeth; MAYO-WILSON, Evan; MCDONALD, Steve; MCGUINNESS, Luke A; STEWART, Lesley A.; THOMAS, James; TRICCO, Andrea C.; WELCH, Vivian; WHITING, Penny; MOHER, David. The PRISMA 2020 statement: an updated guideline for reporting systematic reviews. \textbf{BMJ}, 2021a. DOI: 10.1136/bmj.n71. URL: \url{https://doi.org/10.1136/bmj.n71}. Consulta: 5 out. 2026. Página do trecho não conferida; nível de acesso registrado no fichamento.

\bibitem[Page et al.(2021b)]{SU02} PAGE, Matthew J.; MOHER, David; BOSSUYT, Patrick M.M.; BOUTRON, Isabelle; HOFFMANN, Tammy; MULROW, Cynthia D.; SHAMSEER, Larissa; TETZLAFF, Jennifer; AKL, Elie A.; BRENNAN, Sue E.; CHOU, Roger; GLANVILLE, Julie; GRIMSHAW, Jeremy M.; HRÓBJARTSSON, Asbjørn; LALU, Manoj M.; LI, Tianjing; LODER, Elizabeth; MAYO-WILSON, Evan; MCDONALD, Steve; MCGUINNESS, Luke A; STEWART, Lesley A.; THOMAS, James; TRICCO, Andrea C.; WELCH, Vivian; WHITING, Penny; MCKENZIE, Joanne E.. PRISMA 2020 explanation and elaboration: updated guidance and exemplars for reporting systematic reviews.. \textbf{BMJ}, 2021b. DOI: 10.1136/bmj.n160. URL: \url{https://doi.org/10.1136/bmj.n160}. Consulta: 5 out. 2026. Página do trecho não conferida; nível de acesso registrado no fichamento.

\bibitem[Palalas e Wark(2020)]{education_palalas2020} PALALAS, A.; WARK, Norine. The relationship between mobile learning and self-regulated learning: A systematic review. \textbf{Australasian Journal of Educational Technology}, 2020. DOI: 10.14742/ajet.5650. URL: \url{https://doi.org/10.14742/ajet.5650}. Consulta: 5 out. 2026. Localizador utilizado: p. 156, 164–165. Leitura: \url{https://ajet.org.au/index.php/AJET/article/download/5650/1665}.

\bibitem[Panadero(2017)]{SU10} PANADERO, Ernesto. A Review of Self-regulated Learning: Six Models and Four Directions for Research. \textbf{Frontiers in Psychology}, 2017. DOI: 10.3389/fpsyg.2017.00422. URL: \url{https://doi.org/10.3389/fpsyg.2017.00422}. Consulta: 5 out. 2026. Localizador utilizado: p. 1, 3. Leitura: \url{https://www.frontiersin.org/journals/psychology/articles/10.3389/fpsyg.2017.00422/pdf}.

\bibitem[Patzak et al.(2025)]{education_patzak2025} PATZAK, Alexandra; ZHANG, Xiaorong; VYTASEK, Jovita. Boosting productivity and wellbeing through time management: evidence-based strategies for higher education and workforce development. \textbf{Frontiers in Education}, 2025. DOI: 10.3389/feduc.2025.1623228. URL: \url{https://doi.org/10.3389/feduc.2025.1623228}. Consulta: 5 out. 2026. Página do trecho não conferida; nível de acesso registrado no fichamento.

\bibitem[Prasse et al.(2024)]{education_prasse2024} PRASSE, Doreen; WEBB, Mary; DESCHÊNES, Michelle; PARENT, Séverine; AESCHLIMANN, Franziska; GODA, Yoshiko; YAMADA, Masanori; RAYNAULT, Audrey. Challenges in Promoting Self-Regulated Learning in Technology Supported Learning Environments: An Umbrella Review of Systematic Reviews and Meta-Analyses. \textbf{Technology, Knowledge and Learning}, 2024. DOI: 10.1007/s10758-024-09772-z. URL: \url{https://doi.org/10.1007/s10758-024-09772-z}. Consulta: 5 out. 2026. Localizador utilizado: p. 1809, 1825–1827. Leitura: \url{https://link.springer.com/content/pdf/10.1007/s10758-024-09772-z.pdf}.

\bibitem[Punchoojit e Hongwarittorrn(2017)]{punchoojit2017} PUNCHOOJIT, Lumpapun; HONGWARITTORRN, Nuttanont. Usability Studies on Mobile User Interface Design Patterns: A Systematic Literature Review. \textbf{Advances in Human-Computer Interaction}, 2017. DOI: 10.1155/2017/6787504. URL: \url{https://doi.org/10.1155/2017/6787504}. Consulta: 5 out. 2026. Página do trecho não conferida; nível de acesso registrado no fichamento.

\bibitem[Rethlefsen et al.(2021)]{SU03} RETHLEFSEN, Melissa L.; KIRTLEY, Shona; WAFFENSCHMIDT, Siw; AYALA, Ana Patricia; MOHER, David; PAGE, Matthew J.; KOFFEL, Jonathan; GROUP, PRISMA-S. PRISMA-S: an extension to the PRISMA Statement for Reporting Literature Searches in Systematic Reviews.. \textbf{Systematic Reviews}, 2021. DOI: 10.1186/s13643-020-01542-z. URL: \url{https://doi.org/10.1186/s13643-020-01542-z}. Consulta: 5 out. 2026. Página do trecho não conferida; nível de acesso registrado no fichamento.

\bibitem[Santos et al.(2022)]{eng12} SANTOS, Daniel Amador dos; ALMEIDA, E. D. de; AHMED, Iftekhar. Investigating replication challenges through multiple replications of an experiment. \textbf{Inf. Softw. Technol.}, 2022. DOI: 10.1016/j.infsof.2022.106870. URL: \url{https://doi.org/10.1016/j.infsof.2022.106870}. Consulta: 5 out. 2026. Página do trecho não conferida; nível de acesso registrado no fichamento.

\bibitem[Shahin et al.(2017)]{eng14} SHAHIN, Mojtaba; BABAR, Muhammad Ali; ZHU, Li-Ming. Continuous Integration, Delivery and Deployment: A Systematic Review on Approaches, Tools, Challenges and Practices. \textbf{IEEE Access}, 2017. DOI: 10.1109/ACCESS.2017.2685629. URL: \url{https://doi.org/10.1109/ACCESS.2017.2685629}. Consulta: 5 out. 2026. Página do trecho não conferida; nível de acesso registrado no fichamento.

\bibitem[Shepperd et al.(2018)]{eng13} SHEPPERD, M.; AJIENKA, N.; COUNSELL, S.. The role and value of replication in empirical software engineering results. \textbf{Inf. Softw. Technol.}, 2018. DOI: 10.1016/j.infsof.2018.01.006. URL: \url{https://doi.org/10.1016/j.infsof.2018.01.006}. Consulta: 5 out. 2026. Página do trecho não conferida; nível de acesso registrado no fichamento.

\bibitem[Snyder(2019)]{SU11} SNYDER, Hannah. Literature review as a research methodology: An overview and guidelines. \textbf{Journal of Business Research}, 2019. DOI: 10.1016/j.jbusres.2019.07.039. URL: \url{https://doi.org/10.1016/j.jbusres.2019.07.039}. Consulta: 5 out. 2026. Localizador utilizado: p. 336–337. Leitura: \url{https://fenix.iseg.ulisboa.pt/downloadFile/281608120804966/Guidelines%20Snyder%202019.pdf}.

\bibitem[Soares et al.(2022)]{eng15} SOARES, Eliezio; SIZÍLIO, Gustavo; SANTOS, Jadson; COSTA, Daniel Alencar da; KULESZA, U.. The effects of continuous integration on software development: a systematic literature review. \textbf{Empirical Software Engineering}, 2022. DOI: 10.1007/s10664-021-10114-1. URL: \url{https://doi.org/10.1007/s10664-021-10114-1}. Consulta: 5 out. 2026. Página do trecho não conferida; nível de acesso registrado no fichamento.

\bibitem[Squarcina et al.(2021)]{OS01} SQUARCINA, Marco; CALZAVARA, Stefano; MAFFEI, Matteo. The Remote on the Local: Exacerbating Web Attacks Via Service Workers Caches. \textbf{2021 IEEE Security and Privacy Workshops (SPW)}, 2021. DOI: 10.1109/SPW53761.2021.00062. URL: \url{https://doi.org/10.1109/SPW53761.2021.00062}. Consulta: 5 out. 2026. Página do trecho não conferida; nível de acesso registrado no fichamento.

\bibitem[Subramani et al.(2022)]{OS02} SUBRAMANI, Karthika; JUECKSTOCK, Jordan; KAPRAVELOS, Alexandros; PERDISCI, Roberto. SoK: Workerounds - Categorizing Service Worker Attacks and Mitigations. \textbf{2022 IEEE 7th European Symposium on Security and Privacy (EuroS\&P)}, 2022. DOI: 10.1109/EuroSP53844.2022.00041. URL: \url{https://doi.org/10.1109/EuroSP53844.2022.00041}. Consulta: 5 out. 2026. Página do trecho não conferida; nível de acesso registrado no fichamento.

\bibitem[Sun et al.(2023)]{OS17} SUN, Ruoxi; XUE, Minhui; TYSON, Gareth; WANG, Shuang; CAMTEPE, Seyit; NEPAL, Surya. Not Seen, Not Heard in the Digital World! Measuring Privacy Practices in Children’s Apps. \textbf{ACM Web Conference 2023, pp.2166–2177}, 2023. DOI: 10.1145/3543507.3583327. URL: \url{https://doi.org/10.1145/3543507.3583327}. Consulta: 5 out. 2026. Página do trecho não conferida; nível de acesso registrado no fichamento.

\bibitem[Taghavi-Nejad et al.(2024)]{education_taghavi2024} TAGHAVI-NEJAD, Fatemeh Sadat; FALLAH, Nasser; GASKAREE, Behruz Lotfi. Mindfulness and Procrastination Among University EFL Learners: The Role of Attention Control and Self-Regulated Learning. \textbf{Psychological Reports}, 2024. DOI: 10.1177/00332941241287423. URL: \url{https://doi.org/10.1177/00332941241287423}. Consulta: 5 out. 2026. Página do trecho não conferida; nível de acesso registrado no fichamento.

\bibitem[Tao et al.(2025)]{education_tao2025} TAO, Xue; HANIF, H.; WANG, Lie-Qin. The effects of self-regulated learning strategies on academic procrastination and academic success among college EFL students in China. \textbf{Frontiers in Psychology}, 2025. DOI: 10.3389/fpsyg.2025.1562980. URL: \url{https://doi.org/10.3389/fpsyg.2025.1562980}. Consulta: 5 out. 2026. Página do trecho não conferida; nível de acesso registrado no fichamento.

\bibitem[Tramontana et al.(2018)]{eng03} TRAMONTANA, Porfirio; AMALFITANO, Domenico; AMATUCCI, Nicola; FASOLINO, A.. Automated functional testing of mobile applications: a systematic mapping study. \textbf{Software Quality Journal}, 2018. DOI: 10.1007/s11219-018-9418-6. URL: \url{https://doi.org/10.1007/s11219-018-9418-6}. Consulta: 5 out. 2026. Página do trecho não conferida; nível de acesso registrado no fichamento.

\bibitem[W3C(2024)]{wcag} W3C. Web Content Accessibility Guidelines 2.2, Recommendation de 12 dezembro 2024. \url{https://www.w3.org/TR/2024/REC-WCAG22-20241212/}. Acesso em 5 out. 2026. Registro técnico, fora da contagem de fontes científicas.

\bibitem[Walters e Wilder(2023)]{SU06} WALTERS, William H.; WILDER, Esther Isabelle. Fabrication and errors in the bibliographic citations generated by ChatGPT. \textbf{Scientific Reports, 13, article 14045}, 2023. DOI: 10.1038/s41598-023-41032-5. URL: \url{https://doi.org/10.1038/s41598-023-41032-5}. Consulta: 5 out. 2026. Página do trecho não conferida; nível de acesso registrado no fichamento.

\bibitem[Weigand e Johannesson(2023)]{eng07} WEIGAND, H.; JOHANNESSON, P.. How to Identify your Design Science Research Artifact. \textbf{2023 IEEE 25th Conference on Business Informatics (CBI)}, 2023. DOI: 10.1109/CBI58679.2023.10187511. URL: \url{https://doi.org/10.1109/CBI58679.2023.10187511}. Consulta: 5 out. 2026. Página do trecho não conferida; nível de acesso registrado no fichamento.

\bibitem[Werfhorst et al.(2022)]{werfhorst2022} WERFHORST, H. G. van de; KESSENICH, Emma; GEVEN, Sara. The digital divide in online education: Inequality in digital readiness of students and schools. \textbf{Computers and Education Open}, 2022. DOI: 10.1016/j.caeo.2022.100100. URL: \url{https://doi.org/10.1016/j.caeo.2022.100100}. Consulta: 5 out. 2026. Página do trecho não conferida; nível de acesso registrado no fichamento.

\bibitem[Wolters e Brady(2021)]{education_wolters2021} WOLTERS, Christopher A.; BRADY, Anna C.. College Students’ Time Management: a Self-Regulated Learning Perspective. \textbf{Educational Psychology Review}, 2021. DOI: 10.1007/s10648-020-09519-z. URL: \url{https://doi.org/10.1007/s10648-020-09519-z}. Consulta: 5 out. 2026. Página do trecho não conferida; nível de acesso registrado no fichamento.

\bibitem[Wolters et al.(2017)]{education_wolters2017} WOLTERS, C.; WON, Sungjun; HUSSAIN, M.. Examining the relations of time management and procrastination within a model of self-regulated learning. \textbf{Metacognition and Learning}, 2017. DOI: 10.1007/s11409-017-9174-1. URL: \url{https://doi.org/10.1007/s11409-017-9174-1}. Consulta: 5 out. 2026. Página do trecho não conferida; nível de acesso registrado no fichamento.

\bibitem[Wong et al.(2026)]{education_wong2026} WONG, Jacqueline; KHALIL, Mohammad; SUSCHEVSKIY, Vsevolod; BAARS, Martine; KONING, Bjorn de; PAAS, Fred. Student engagement profiles in a mobile app: Links to self-regulated learning and performance. \textbf{Educational Technology Research and Development}, 2026. DOI: 10.1007/s11423-026-10586-2. URL: \url{https://doi.org/10.1007/s11423-026-10586-2}. Consulta: 5 out. 2026. Página do trecho não conferida; nível de acesso registrado no fichamento.

\bibitem[Xu et al.(2022)]{education_xu2022} XU, Zhi-Hong; ZHAO, Yingying; ZHANG, Bingsheng; LIEW, J.; KOGUT, Ashlynn. A meta-analysis of the efficacy of self-regulated learning interventions on academic achievement in online and blended environments in K-12 and higher education. \textbf{Behaviour \& Information Technology}, 2022. DOI: 10.1080/0144929X.2022.2151935. URL: \url{https://doi.org/10.1080/0144929X.2022.2151935}. Consulta: 5 out. 2026. Página do trecho não conferida; nível de acesso registrado no fichamento.

\bibitem[Zein et al.(2023)]{eng04} ZEIN, Samer; SALLEH, Norsaremah; GRUNDY, John C.. Systematic reviews in mobile app software engineering: A tertiary study. \textbf{Inf. Softw. Technol.}, 2023. DOI: 10.1016/j.infsof.2023.107323. URL: \url{https://doi.org/10.1016/j.infsof.2023.107323}. Consulta: 5 out. 2026. Página do trecho não conferida; nível de acesso registrado no fichamento.

\bibitem[Zein et al.(2016)]{eng05} ZEIN, Samer; SALLEH, Norsaremah; GRUNDY, J.. A systematic mapping study of mobile application testing techniques. \textbf{J. Syst. Softw.}, 2016. DOI: 10.1016/j.jss.2016.03.065. URL: \url{https://doi.org/10.1016/j.jss.2016.03.065}. Consulta: 5 out. 2026. Página do trecho não conferida; nível de acesso registrado no fichamento.

\bibitem[Zhao et al.(2022)]{OS15} ZHAO, Jun; DURON, Blanche; WANG, Ge. KOALA Hero: Inform Children of Privacy Risks of Mobile Apps. \textbf{Interaction Design and Children, 2022}, 2022. DOI: 10.1145/3501712.3535278. URL: \url{https://doi.org/10.1145/3501712.3535278}. Consulta: 5 out. 2026. Página do trecho não conferida; nível de acesso registrado no fichamento.

\bibitem[Zheng et al.(2016)]{education_zheng2016} ZHENG, Lanqin; LI, Xin; CHEN, Fengying. Effects of a mobile self-regulated learning approach on students’ learning achievements and self-regulated learning skills. \textbf{Innovations in Education and Teaching International}, 2016. DOI: 10.1080/14703297.2016.1259080. URL: \url{https://doi.org/10.1080/14703297.2016.1259080}. Consulta: 5 out. 2026. Página do trecho não conferida; nível de acesso registrado no fichamento.

\end{thebibliography}
\newpage\section*{Apêndices}\addcontentsline{toc}{section}{Apêndices}
\textbf{Apêndice A — roteiro técnico elaborado para o projeto.} Criar uma tarefa sintética; concluir; recarregar; exportar backup; limpar dados em ambiente de teste; importar; conferir valores. Tentar backup inválido e confirmar rejeição sem perda dos dados anteriores. Instalar cache online; fechar a aplicação; bloquear rede; reabrir e registrar resultado. Nunca executar limpeza em dados reais sem backup conferido; o roteiro detalhado é o protocolo do projeto \citep{projeto}.
\section*{Anexos}\addcontentsline{toc}{section}{Anexos}
\textbf{Anexo A — documento externo de referência.} Modelo ``01-ENTREGA-FINAL-PROJETO-INTEGRADOR-9-1-.docx'', fornecido para esta demonstração. O original não é reproduzido; foi usado para conferir a estrutura desta entrega. Não foram incluídos assinaturas, certificados ou documentos institucionais não fornecidos, conforme o registro de requisitos do projeto \citep{projeto}.
\end{document}